\noindent\textbf{(a)} $x^2\cos y+x^3y-y^2/2=C$.

\noindent\textbf{(b)} $x^3+3x^2y+3xy^2-3y=4$. Su rama inicial es\[y=\frac{3-3x^2+\sqrt{9-18x^2-3x^4+48x}}{6x},\quad 0<x<\alpha,\]donde $\alpha\in(1,2)$ es la única raíz positiva de $x^4+6x^2-16x-3=0$.

\noindent\textbf{(c)} $y^2\sin x-x^3y-x^2+y\log y-y=0$, en la rama regular que pasa por $(0,e)$ y permanece en $y>0$.

\noindent\textbf{(d)} $x^2/y^2+xy=C$ para $y\ne0$, y además $y=0$. Por el origen hay además prolongaciones $C^1$ de las ramas $C>0$, con pendiente $\pm1/\sqrt C$; la misma constante y el mismo signo de pendiente deben mantenerse a ambos lados.

\noindent\textbf{(e)} $x^2/y+xy+y^2=C$ para $y\ne0$, y además $y=0$. Por el origen, cada $C\ne0$ tiene una rama analítica con $y=x^2/C+O(x^5)$ y $y^{\prime}(0)=0$. Pueden empalmarse a cada lado esas ramas con constantes independientes o con la rama cero, conservando clase $C^1$. El nivel $C=0$ no aporta una rama no nula $C^1$ por el origen.
